\documentclass[11pt,a4paper]{article}
\usepackage[utf8]{inputenc}
\usepackage[T1]{fontenc}
\usepackage[spanish]{babel}
\usepackage{lmodern}
\usepackage{geometry}
\usepackage{booktabs}
\usepackage{array}
\usepackage{microtype}
\usepackage{xcolor}
\usepackage{xurl}
\usepackage[hidelinks]{hyperref}
\geometry{margin=2.3cm, bottom=2.5cm}
\setlength{\parskip}{0.6em}
\setlength{\parindent}{0pt}
\renewcommand{\arraystretch}{1.25}

\hypersetup{
  pdftitle={Detección de lavado de dinero en remesas - Reporte Ejecutivo del Proyecto 2},
  pdfauthor={Jonathan Díaz, David Dominguez, Gabriel Bran}
}

\newcommand{\metric}[1]{\texttt{#1}}

\begin{document}

\begin{center}
{\LARGE\bfseries Detección de lavado de dinero en remesas\par}
\vspace{0.3em}
{\large Reporte ejecutivo | CC3092 Deep Learning, Proyecto 2\par}
\vspace{0.6em}
{\normalsize Jonathan Díaz (23837) \quad David Dominguez (23712) \quad Gabriel Bran (23590)\par}
{\small Universidad del Valle de Guatemala \quad|\quad Septiembre de 2026\par}
\end{center}

\vspace{0.6em}
\noindent\textbf{Decisión clave:} el prototipo sirve para priorizar qué remitentes investigar primero, pero no supera de forma clara a un modelo supervisado simple en la capacidad real de revisión (el 1\,\% de alertas). No justifica una operación productiva sin validación temporal y con datos reales.

\section*{Resumen}
Construimos un sistema de dos etapas para detectar patrones de lavado de dinero en el comportamiento de un remitente a lo largo del tiempo. La primera etapa aprende cómo se ve el comportamiento normal y mide qué tan raro es un remitente; la segunda usa casos etiquetados de lavado para decidir si un remitente debe ser alertado, y explica su decisión señalando qué transacciones pesaron más. Se entrenó con IBM AML, un conjunto sintético de 5,078,345 transacciones y 496,995 remitentes, y se evaluó sobre 9,002 remitentes de prueba (245 con lavado etiquetado).

El resultado es matizado y lo reportamos tal como salió. La señal de la primera etapa mejora el ordenamiento global del riesgo (ROC-AUC de 0.858 frente a 0.830 del modelo base, en las tres semillas). Sin embargo, en la métrica que más se acerca a la operación, la precisión con el 1\,\% de alertas, el modelo base es mejor (0.785 frente a 0.662), y en AUPRC la diferencia es de 0.005, dentro del ruido. Además, el sistema falló en un caso con una ráfaga de montos extremos. Por eso lo presentamos como una arquitectura viable y un instrumento de evaluación, no como un sustituto de los controles vigentes.

\section{Problema de negocio y contexto regulatorio}
Guatemala recibe más de USD 20 mil millones anuales en remesas, cerca del 19\,\% del PIB según el planteamiento del proyecto. Ese volumen convierte al corredor Guatemala--Estados Unidos en un canal atractivo para mover fondos ilícitos entre miles de envíos pequeños y aparentemente ordinarios. Las empresas de remesas y los bancos receptores deben detectar patrones sospechosos, pero los sistemas basados en reglas generan entre 95\,\% y 99\,\% de falsos positivos, de modo que el equipo de cumplimiento revisa cientos de alertas para hallar un caso real. Lo que revela lavado no es una transacción aislada sino el patrón en el tiempo: frecuencia inusual, montos fragmentados, destinos que cambian de golpe. El aprendizaje automático supervisado ya se ha usado para esto \cite{jullum2020}, y la literatura señala la explicabilidad como requisito para entornos regulados \cite{kute2021}.

El marco legal guatemalteco cambió recientemente. Según la información pública disponible, el Congreso aprobó en junio de 2026 el Decreto 15-2026, Ley Integral para la Prevención y Represión del Lavado de Dinero u Otros Activos y del Financiamiento del Terrorismo, que unifica y deroga la Ley contra el Lavado de Dinero (Decreto 67-2001) y la Ley para Prevenir y Reprimir el Financiamiento del Terrorismo (Decreto 58-2005) \cite{infobae2026,pirani2026}. La nueva ley entra en vigor tres meses después de su publicación (17 de junio de 2026), es decir, a mediados de septiembre de 2026. Incluye entre los sujetos obligados a las actividades de transferencia de fondos, exige administración de riesgo con enfoque basado en riesgo, manual de prevención, reporte de operaciones sospechosas y auditoría anual de efectividad, y la supervisión recae en la Superintendencia de Bancos a través de la Intendencia de Verificación Especial \cite{pirani2026}. La ley responde a las Recomendaciones del GAFI, evaluadas en la región por GAFILAT \cite{fatf2012,gafilat2016}, por lo que los países vecinos comparten el mismo estándar. Del lado estadounidense, las remesadoras son sujetos de la regulación de FinCEN; el umbral de USD 10,000 que usamos como variable refleja el límite de reporte de operaciones en efectivo y el riesgo de \emph{structuring}, es decir, fraccionar montos para evitarlo.

Dos advertencias. Primero, no citamos artículos concretos de la nueva ley porque no pudimos consultar su texto oficial; la información proviene de reportes públicos y debe contrastarse con la publicación oficial antes de cualquier uso legal. Segundo, el modelo no es una prueba legal: su puntaje es una señal de priorización que requiere revisión humana y evidencia adicional. Por eso el umbral no se eligió por una métrica académica sino por la capacidad de revisión del equipo, asumida en 1\,\% de remitentes alertados.

\section{Datos y decisiones de diseño}
\textbf{Fuente de datos.} El enunciado propone PaySim como conjunto principal. Lo cargamos y lo diagnosticamos: cada remitente aparece en promedio 1.0015 veces, así que es imposible construir un historial. Usamos IBM AML (variante \texttt{HI-Small}) \cite{altman2023}, con remitentes recurrentes (mediana de 2 transacciones, máximo de 168,672); PaySim quedó como diagnóstico secundario. Esta es una decisión nuestra que se aparta del enunciado y que justificamos con esa evidencia.

\textbf{Secuencias.} Cada transacción se describe con 14 variables: monto (en logaritmo), tiempo desde la transacción anterior, hora del día (codificada de forma cíclica), tipo de pago, si el destino es nuevo, entropía acumulada de destinos (qué tan disperso es el abanico de contrapartes) y monto sobre USD 10,000. Se normaliza con estadísticos calculados solo sobre entrenamiento. Descartamos remitentes con menos de 5 transacciones (74\,\% de todos) y truncamos a las últimas 29, valor tomado del percentil 90 de las longitudes. De 60,000 remitentes de la muestra, 1,626 tienen lavado (2.71\,\%), y se conservaron todos los positivos porque son escasos; la prevalencia poblacional es de 0.68\,\%. La partición es 70/15/15 por remitente, estratificada, con una verificación automática de que ningún remitente aparece en dos particiones: entrenamiento 41,999 (1,138 positivos), validación 8,999 (243) y prueba 9,002 (245).

\textbf{Etiqueta.} Un remitente es positivo si tiene alguna transacción de lavado en todo su historial y no solo en la ventana de 29. Con la versión anterior, 163 remitentes con lavado quedaban etiquetados como normales y contaminaban el aprendizaje de la normalidad. El costo es que en 163 positivos el modelo no ve la transacción ilícita.

\textbf{Etapa A: aprender la normalidad.} Es un autoencoder: un codificador (red recurrente GRU bidireccional con una capa de atención) comprime la secuencia en un vector de 64 números, y un decodificador intenta reconstruirla. Se entrena solo con remitentes normales; cuando aparece uno que no se parece a nada visto, la reconstrucción falla y ese error es el puntaje de anomalía. Una primera versión, que repetía el vector en cada paso, casi no aprendía (el error de validación bajó apenas de 0.81 a 0.74); el decodificador final recibe además la transacción anterior real, y el error bajó a 0.42 (el nivel de predecir siempre la media es 1.0). El umbral de anomalía es el que produce 1\,\% de alertas sobre validación (1.269).

\textbf{Etapa B: aprender del lavado etiquetado.} Reutiliza el codificador de la etapa A (transfer learning): una época con el codificador congelado y luego ajuste fino completo, con tasa de aprendizaje menor para el codificador que para la cabeza nueva, para no borrar lo aprendido. Como los casos de lavado son pocos, la pérdida es \emph{focal loss} implementada a mano, que reduce el peso de los ejemplos fáciles \cite{lin2017}. El error de reconstrucción de la etapa A entra como una variable más de la cabeza clasificadora en lugar de promediar puntajes de escalas distintas. El umbral de alerta también se fija en validación (0.334) y se aplica a prueba: produjo 105 alertas (1.17\,\%).

\textbf{Conexión con el curso.} La GRU (S3) evita el desvanecimiento del gradiente de una RNN simple en secuencias de hasta 29 pasos; el autoencoder con decodificador autoregresivo (S4) es el esquema codificador--decodificador aplicado a detección de fraude; la atención aditiva (S5) da los pesos que se visualizan en el análisis; la focal loss, Adam y el \emph{dropout} (S8) atienden el desbalance y la regularización; y congelar y luego ajustar con tasas distintas (S9) es el protocolo gradual de transferencia. Detrás de todo, el desbalance justifica priorizar AUPRC sobre ROC-AUC.

\section{Resultados de la ablación}
Comparamos cinco configuraciones con tres semillas cada una y reportamos media y desviación estándar. AUPRC resume qué tan bien se ordenan los casos de lavado por encima de los normales; un modelo al azar da 0.027, la prevalencia. La \metric{Precision@1\%} es la fracción de las alertas que eran lavado real, y el \metric{Recall@1\%} la fracción de los casos de lavado que se logró alertar.

\begin{table}[htbp]
\centering
ootnotesize
\setlength{	abcolsep}{4pt}
\caption{Resultados de ablación en prueba (media $\pm$ desviación estándar, 3 semillas)}
\label{tab:ablacion}
\begin{tabular}{lccccc}
\toprule
Configuración & AUPRC & ROC-AUC & P@1\% & R@1\% & F2 \\
\midrule
1. Base supervisada (sin etapa A) & 0.398 $\pm$ 0.002 & 0.830 $\pm$ 0.001 & 0.785 $\pm$ 0.014 & 0.307 $\pm$ 0.008 & 0.350 \\
2. Solo etapa A & 0.060 $\pm$ 0.001 & 0.723 $\pm$ 0.001 & 0.073 $\pm$ 0.007 & 0.027 $\pm$ 0.005 & 0.031 \\
3. Codificador congelado & 0.107 $\pm$ 0.013 & 0.688 $\pm$ 0.015 & 0.301 $\pm$ 0.045 & 0.088 $\pm$ 0.012 & 0.103 \\
4. Ajuste fino & 0.379 $\pm$ 0.016 & 0.832 $\pm$ 0.003 & 0.681 $\pm$ 0.015 & 0.294 $\pm$ 0.022 & 0.332 \\
5. Ajuste fino + reconstrucción & 0.403 $\pm$ 0.012 & 0.858 $\pm$ 0.003 & 0.662 $\pm$ 0.022 & 0.294 $\pm$ 0.018 & 0.331 \\
\bottomrule
\end{tabular}
\end{table}

Qué aporta cada pieza. La etapa A sola (2) detecta unas 2.2 veces mejor que el azar: aprender lo normal no basta para señalar lavado. Congelar el codificador (3) queda muy por debajo, porque pasar de reconstruir a clasificar exige adaptarlo. El ajuste fino sin reconstrucción (4) no supera a la base: el codificador preentrenado, por sí solo, no aporta. Lo que sí aporta es la señal de reconstrucción (5): sube el AUPRC de 0.379 a 0.403 respecto de la configuración 4 y el ROC-AUC de 0.830 a 0.858 respecto de la base, con mejora en las tres semillas. Frente a la base, en cambio, el AUPRC mejora solo 0.005 y no de forma consistente: gana en dos semillas y pierde en una.

En el punto operativo el resultado se invierte. Con 1\,\% de alertas, la base logra precisión de 0.785 frente a 0.662 de la configuración 5, y mayor F2, en las tres semillas. Traducido a la operación: de las 105 alertas del sistema completo, 70 eran lavado etiquetado (66.7\,\%); la base habría entregado una cola más limpia. La conclusión es que este experimento \emph{no demuestra} que las dos etapas superen a la base, y solo respalda una mejora del ordenamiento global. Una comparación con el 95--99\,\% de falsos positivos de los sistemas de reglas tampoco es válida aquí, porque los datos son sintéticos y la prueba está enriquecida (2.7\,\% de positivos).

\section{Análisis de casos}
Los pesos de atención indican cuánta importancia dio el modelo a cada transacción; son una pista de dónde mirar, no una prueba de que esa transacción sea ilícita. Las etiquetas son por remitente y no por transacción. Analizamos cinco casos de prueba; las tipologías se leen sobre los datos reales y no sobre la heurística automática, que etiquetó los cinco como ``fraccionamiento'' por un defecto que describimos más abajo.

\textbf{Caso 0, verdadero positivo (probabilidad 98.2\,\%).} La atención (59\,\%) cae en un tramo final de transferencias ACH de USD 24.4 mil tras 55 horas de inactividad, USD 4.2 mil tras 84 horas y USD 16.1 mil, hacia la propia cuenta. El historial previo mezcla cheques, tarjeta, efectivo y \emph{wire}, y repite montos idénticos (USD 203,402.36, 61,128.98 y 27,815.69, dos veces cada uno). Encaja con \emph{layering}: cambio de canal, escalada de montos y repeticiones exactas. Solo la etapa B lo alerta; la etapa A lo ve por debajo de su umbral.

\textbf{Caso 1, verdadero positivo (97.8\,\%).} La atención se concentra (91\,\%) en el cierre: una ACH de USD 15.9 mil, otra de 3.9 mil a las 23 h, once minutos después, y un cheque de 8.9 mil a la misma contraparte. Antes hay dos pares idénticos de montos de más de un millón de dólares. Su última transacción va al remitente del caso 0, también con lavado: un vínculo entre dos casos que un modelo que mira cada remitente por separado no puede aprovechar, y que justifica explorar enfoques con grafos \cite{cheng2023}.

\textbf{Caso 2, verdadero positivo (95.8\,\%).} Es el único donde alertan ambas etapas. Reparte 29 transacciones entre 15 destinos, 12 de ellos nuevos: repite tres veces un ciclo de cuatro montos y luego encadena 16 transferencias ACH de USD 5.6 a 70.2 mil. Es el patrón de \emph{abanico de destinos}, es decir, dispersión hacia muchas contrapartes.

\textbf{Caso 3, falso positivo (60.5\,\%).} Alertan ambas etapas, pero el remitente figura como normal. La atención se concentra en una ráfaga de transacciones moderadas (USD 1.2 a 3.8 mil, dos de ellas separadas por 15 minutos) que termina en un destino nuevo. No hay montos extremos; una hipótesis es actividad legítima, como pagos agrupados a proveedores, con la misma huella que una ilícita. No podemos verificarla, y un analista tampoco sin datos de identificación del cliente: es el costo típico de un falso positivo.

\textbf{Caso 4, falso negativo (4.0\,\%).} Es la falla más importante. Las primeras 18 transacciones alternan exactamente dos montos (USD 16,615.62 y 18,464.82) hacia tres destinos, una regularidad que el autoencoder reconstruye casi perfecto y por eso parece ``normal''. Después hay una ráfaga de tres ACH por USD 201.9 millones, 29.2 millones y 0.9 millones, de órdenes de magnitud sobre el resto, y el modelo les dio solo 1.4\,\% de la atención. El historial completo (23 transacciones) estaba visible, así que no fue un problema de truncamiento. Una regla simple de monto extremo lo habría capturado, lo que muestra que el modelo debe complementar a las reglas y no reemplazarlas.

\section{Limitaciones}
\textbf{Datos sintéticos.} Las relaciones aprendidas pueden no reflejar carteras reales ni tácticas que se adaptan. \textbf{Partición aleatoria.} No hay validación temporal estricta; en producción se entrena con un periodo y se evalúa con los posteriores. \textbf{Cobertura.} El filtro de 5 transacciones excluye 1,750 de los 3,376 remitentes con lavado (52\,\%): el sistema no ve a los que operan poco. \textbf{Prueba enriquecida.} La muestra tiene 2.7\,\% de positivos frente a 0.68\,\% en la población, por lo que la precisión no se traslada a producción. \textbf{Ruido de etiqueta.} 163 positivos tienen el lavado fuera de la ventana visible. \textbf{Heurística de tipologías.} Marca como ``fraccionamiento'' cualquier monto sobre USD 8,000, y eso ocurre en 1,110 de los 2,000 remitentes del MVP; debe acotarse a montos justo bajo el umbral. \textbf{Longitud máxima.} El percentil 90 se calculó antes del filtro y dio 29, menor que lo que daría sobre los remitentes válidos. \textbf{Una sola semilla en los casos.} Los cinco casos y el MVP usan el modelo de la semilla 0. \textbf{Atención.} No prueba causalidad. \textbf{MVP parcial.} Muestra 2,000 de los 9,002 remitentes de prueba (los 245 positivos y 1,755 negativos al azar), no incluye el falso positivo del caso 3 y no ejecuta el modelo en línea. \textbf{Reproducibilidad.} El notebook corre completo en 17.8 minutos en una GPU T4 de Colab.

\section{Ruta a producción}
Se necesitan datos históricos con investigaciones concluidas y etiquetas con linaje, enriquecidos con identificación del cliente, beneficiario final, canal y geografía, y una validación temporal que respete cuándo se conocía cada dato. El modelo debe entregar una cola priorizada con explicaciones y evidencia navegable; la decisión final es del analista, que documenta el resultado y alimenta un ciclo de aprendizaje gobernado, no un reentrenamiento automático. Deben existir trazas de auditoría de la versión que generó cada alerta, coherentes con el manual de prevención y la auditoría anual que exige la nueva ley. Se recomienda combinarlo con reglas de monto extremo, calibrar el umbral por capacidad y costo de error, monitorear la deriva de datos y del desempeño, y probar primero en un piloto controlado.

\section{Conclusión}
El proyecto entrega un modelo secuencial evaluado con ablación, métricas pertinentes al desbalance y casos que muestran aciertos y fallos. La señal de reconstrucción mejora el ordenamiento global, pero la base supervisada conserva ventaja en el uno por ciento superior de revisión, y no queda demostrado que las dos etapas la superen. La decisión responsable es tratar el prototipo como instrumento de aprendizaje y piloto analítico, con reglas simples como red de seguridad, y corregir los defectos señalados antes de cualquier uso operativo.

\section*{Enlaces del proyecto}
\textbf{Repositorio público:} \url{https://github.com/Branuvg/labs_deep}

\textbf{MVP público:} \url{https://labsdeep-proyecto2.streamlit.app/}

\begin{thebibliography}{99}
\small
\bibitem{jullum2020} Jullum, M., Løland, A., Huseby, R. B., Ånonsen, G., \& Lorentzen, J. (2020). Detecting money laundering transactions with machine learning. \textit{Journal of Money Laundering Control, 23}(1), 173--186. \url{https://doi.org/10.1108/JMLC-07-2019-0055}
\bibitem{kute2021} Kute, D. V., Pradhan, B., Shukla, N., \& Alamri, A. (2021). Deep learning and explainable artificial intelligence techniques applied for detecting money laundering - A critical review. \textit{IEEE Access, 9}, 82300--82317. \url{https://doi.org/10.1109/ACCESS.2021.3086230}
\bibitem{cheng2023} Cheng, D., Ye, Y., Xiang, S., Ma, Z., Zhang, Y., \& Jiang, C. (2023). Anti-Money Laundering by Group-Aware Deep Graph Learning. \textit{IEEE Transactions on Knowledge and Data Engineering, 35}(12), 12444--12457. \url{https://doi.org/10.1109/TKDE.2023.3272396}
\bibitem{altman2023} Altman, E., Blanuša, J., von Niederhäusern, L., Egressy, B., Anghel, A., \& Atasu, K. (2023). Realistic synthetic financial transactions for anti-money laundering models. \textit{Advances in Neural Information Processing Systems 36 (NeurIPS), Datasets and Benchmarks Track}. \url{https://arxiv.org/abs/2306.16424}
\bibitem{lin2017} Lin, T.-Y., Goyal, P., Girshick, R., He, K., \& Dollár, P. (2017). Focal loss for dense object detection. \textit{IEEE International Conference on Computer Vision (ICCV)}. \url{https://doi.org/10.1109/ICCV.2017.324}
\bibitem{fatf2012} Financial Action Task Force. (2012-2023). \textit{International Standards on Combating Money Laundering and the Financing of Terrorism \& Proliferation: The FATF Recommendations}. \url{https://www.fatf-gafi.org/en/publications/Fatfrecommendations/Fatf-recommendations.html}
\bibitem{gafilat2016} GAFILAT. (2016). \textit{Informe de evaluación mutua de la República de Guatemala}. \url{https://www.gafilat.org/index.php/es/biblioteca-virtual/miembros/guatemala/evaluaciones-mutuas-11/3374-informe-de-evaluacion-mutua-de-guatemala/file}
\bibitem{infobae2026} Infobae. (3 de junio de 2026). Congreso de Guatemala aprueba Ley Antilavado de Dinero. \url{https://www.infobae.com/guatemala/2026/06/03/congreso-de-guatemala-aprueba-ley-antilavado-de-dinero-en-favor-de-la-transparencia-financiera/}
\bibitem{pirani2026} Pirani Risk. (2026). Ley 15-2026: nueva ley antilavado en Guatemala, qué es y a quién aplica. \url{https://www.piranirisk.com/es/hub-regulatorio/ley-15-2026-antilavado-guatemala-que-es-a-quien-aplica}
\end{thebibliography}

\section*{Declaración de uso de IA}
Declaración de uso de IA: Se utilizó asistencia de IA generativa para apoyar la redacción, organización y revisión de claridad de este reporte. El contenido factual se limitó al dossier proporcionado para el Proyecto 2; las métricas, limitaciones y casos no fueron generados como resultados experimentales. La responsabilidad de verificar las fuentes, revisar el PDF y aprobar el contenido final corresponde al autor del trabajo.

\end{document}
